\section{Просодические характеристики речевого сигнала}

\subsection{Основные понятия просодики (задание 2.1)}

Просодика описывает две основные составляющие речевого сигнала, которые
не связаны с фонемным составом слова, но передают дополнительный смысл
высказывания: \emph{ударение} и \emph{интонацию}~\cite{lobanov}.

Интонация присуща естественной речи: без неё речь звучит монотонно и
неестественно и трудно воспринимается, поскольку вместе с интонацией
говорящий передаёт эмоциональную и смысловую окраску высказывания.
Ударение позволяет различать слова и понимать смысл сообщений; кроме
словесного ударения в просодике рассматриваются фразовое и логическое
ударения, которые делают речь более выразительной и участвуют в передаче
интонации.

\subsection{Акцентная группа}

Акцентная группа (АГ) делится на три составляющие --- элементы акцентной
группы (ЭАГ):
\begin{itemize}
	\item \emph{предъядро} --- все звуки слева от главного ударного гласного;
	\item \emph{ядро} --- главная ударная гласная слова;
	\item \emph{заядро} --- все звуки справа от ядра.
\end{itemize}
Предъядро и заядро могут содержать различное количество звуков вплоть до
их отсутствия. Например, во фразе <<Ты поедешь сегодня? --- Да!>> слово
<<Да!>> одновременно является синтагмой и акцентной группой, поэтому
единственный гласный звук приходится делить на три части, чтобы выделить
ядро, предъядро и заядро.

\subsection{Интонационные типы конечной акцентной группы}

Интонационный тип фразы определяется положением мелодического максимума
$F_0$ относительно главного ударного гласного (ядра):
\begin{itemize}
	\item сдвиг максимума \emph{влево} от центра ядра соответствует
	интонации завершения или утверждения (на письме --- точка <<.>>);
	\item положение максимума \emph{в центре} гласного соответствует
	интонации побуждения или восклицания (<<!>>);
	\item сдвиг максимума \emph{вправо} соответствует интонации
	незавершённости (<<,>>, <<->>) или вопроса (<<?>>).
\end{itemize}
Другие параметры мелодической кривой --- значение максимума и его ширина
--- существенны при передаче эмоциональных оттенков речи.

\subsection{Используемое программное обеспечение}

Запись речевого сигнала выполнялась в программе Audacity 4.0.1
(частота дискретизации 16~кГц, моно, 16 бит), которая, как и
Sonic Foundry Sound Forge, позволяет записывать сигнал с микрофона,
просматривать осциллограмму, выделять, вырезать и вставлять фрагменты,
изменять громкость, а также сохранять результат в файл формата
\texttt{*.wav}.

Обработка и анализ сигналов выполнялись на языке Python с использованием
библиотек NumPy, SciPy, soundfile и Matplotlib:
\begin{itemize}
	\item сегментация речевого сигнала по кратковременной энергии
	(кадры 10~мс, порог относительно максимума, объединение близких
	сегментов);
	\item оценивание частоты основного тона $F_0$ методом
	автокорреляции в диапазоне 60--400~Гц с контролем силы периодичности
	и подавлением октавных выбросов;
	\item аналог \texttt{Effect -> Pitch -> Bend} --- плавное изменение
	$F_0$ по заданному контуру с сохранением длительности (ресемплинг +
	фазовый вокодер, блочная обработка с перекрытием блоков);
	\item аналог \texttt{Effects -> Pitch -> Shift} --- постоянный сдвиг
	тона всего сигнала;
	\item аналог \texttt{Process -> Fade -> Graphic} --- плавное изменение
	амплитуды выделенных слов;
	\item аналог \texttt{Process -> Time Stretch} --- изменение
	длительности ударного слова без изменения тона.
\end{itemize}
